\section{Reconstruction algebra and bounded residual decoding}
\label{app:reconstruction-mechanism}

This appendix analyzes fixed reference maps and their use by a specified
decoder. It establishes normalized-kernel identities, conditions for equality
or inequality of feature spans, and a bound on fitted residual corrections.
None of these statements assumes a particular time-series distribution or
guarantees a forecasting improvement. Unless another subscript is displayed,
vector norms are Euclidean norms. The numerical examples can be evaluated
with \path{extensions/reconstruction_mechanism/algebra.py}; their reference
banks are constructed independently of forecasting outcomes.

\subsection{Normalized kernel identities}

Let $q_1,\ldots,q_M\in\R^d$ be reference histories, and let
$Q\in\R^{M\times d}$ contain their transposes as rows. Let
$P\in\R^{d\times r}$ have orthonormal columns, with $1\leq r\leq d$,
and let $W\in\R^{r\times r}$ be a symmetric positive-definite precision
matrix for the projected coordinates. For a column-vector query
$z\in\R^d$, define the Gaussian kernel vector $k(z)\in\R^M$ by
\[
 k_i(z)=\exp\!\left[-\frac{
 (P^\top(z-q_i))^\top W(P^\top(z-q_i))}{2r}\right],
 \qquad i=1,\ldots,M.
\]
Let $K\in\R^{M\times M}$ be the reference Gram matrix, with
$K_{ij}=k_i(q_j)$. For $\lambda>0$, set
$C=(K+\lambda I_M)^{-1}$, where $I_M$ is the $M\times M$ identity
matrix, and let $\mathbf1_M$ be the length-$M$ vector of ones.
The matrix $K$ is symmetric positive semidefinite, so $C$ exists and is
symmetric positive definite. The KAHM and matched positive-kernel
reconstructions are
\begin{equation}
 A_K(z)=\frac{Q^\top Ck(z)}{\mathbf1_M^\top Ck(z)},\qquad
 A_R(z)=\frac{Q^\top k(z)}{\mathbf1_M^\top k(z)}.
 \label{eq:mechanism-normalized-maps}
\end{equation}
The first expression is defined wherever its signed denominator is nonzero.
The second denominator is positive for every finite query. The two maps
share the references and projected Gaussian kernel. Multiplying $k(z)$
by a common positive scalar cancels in either ratio, which justifies
subtracting the largest log-kernel value before query evaluation.

\paragraph{Target-transfer identity.}
For a query satisfying $s(z)=\mathbf1_M^\top Ck(z)\ne0$, define the
normalized weight vector $\omega(z)=Ck(z)/s(z)$. Its entries sum to one.
For any reference target vector $y\in\R^M$,
\[
 \omega(z)^\top y=
 \frac{k(z)^\top(K+\lambda I_M)^{-1}y}
 {k(z)^\top(K+\lambda I_M)^{-1}\mathbf1_M}.
\]
Thus target transfer is the ratio of two dual kernel-ridge predictions
with regularization matrix $K+\lambda I_M$ and no additional intercept:
one fitted to $y$, and one fitted to the constant target $\mathbf1_M$.
For an affine scalar function $f(q)=q^\top b+c$, with
$b\in\R^d$ and $c\in\R$, let $f(Q)=Qb+c\mathbf1_M$ be its reference
values. Then
\[
 \omega(z)^\top f(Q)-f(z)=\{A_K(z)-z\}^\top b.
\]
\begin{proof}[Proof of Proposition~\ref{prop:mechanism-krr}]
The definition gives
$\mathbf1_M^\top\omega(z)=s(z)/s(z)=1$. Because $C=C^\top$,
$\omega(z)^\top y=k(z)^\top Cy/s(z)$ and
$s(z)=k(z)^\top C\mathbf1_M$, proving the ratio identity.
For the affine function,
\[
 \omega(z)^\top f(Q)
 =\omega(z)^\top Qb+c\omega(z)^\top\mathbf1_M
 =A_K(z)^\top b+c.
\]
Subtracting $z^\top b+c$ proves the final identity.
\end{proof}

This identifies direct target transfer as a normalized kernel-ridge
construction. The ratio generally differs from the unnormalized
kernel-ridge predictor because its denominator depends on the query.
Similarly, $A_K(z)^\top b$ is normalized kernel ridge applied to the
restricted reference target $Qb$. For any reconstruction $A$, the identity
$A(z)+\{z-A(z)\}=z$ also implies a linear dependence among the raw input,
reconstruction, and residual when all three are included as unrestricted
linear features. These identities constrain claims about the resulting
function class while allowing useful parameterizations for a restricted
or regularized decoder.

\subsection{When reconstruction departures have the same span}

Let $\mu\in\R^d$ be the reference center and define the affine
reconstruction $A_A(z)=\mu+PP^\top(z-\mu)$. Fix a common scalar locality
gate $g(z)$, and define column-vector departures
\[
 d_K(z)=g(z)\{A_K(z)-A_A(z)\},\qquad
 d_R(z)=g(z)\{A_R(z)-A_A(z)\}.
\]
A feature span means the set of scalar functions obtained by finite linear
combinations of its coordinate functions. Let $\mathcal S_0$ be the span
of the gate $g(z)$ and the coordinates of $g(z)P^\top(z-\mu)$.
It contains every coordinate of $g(z)A_A(z)$.

\paragraph{Equal augmented feature spans.}
Suppose $K\mathbf1_M=\kappa\mathbf1_M$ for a scalar $\kappa$ and
$C$ leaves the augmented reference-coordinate space
$\mathcal V=\operatorname{span}\{\mathbf1_M,Q_{:1},\ldots,Q_{:d}\}
\subseteq\R^M$ invariant. Here $Q_{:j}$ denotes column $j$ of $Q$.
Then there exist $B\in\R^{d\times d}$ and $c\in\R^d$ such that
\[
 C\mathbf1_M=a\mathbf1_M,\qquad
 CQ=QB+\mathbf1_M c^\top,\qquad a=(\kappa+\lambda)^{-1}>0.
\]
For these coefficients,
\[
 A_K=\frac{B^\top A_R+c}{a},\qquad
 d_K=\frac{B^\top}{a}d_R+
 \left(\frac{B^\top}{a}-I_d\right)gA_A+\frac{gc}{a}.
\]
The span obtained by adjoining the coordinates of $d_K$ to
$\mathcal S_0$ is the same as that obtained by adjoining the coordinates
of $d_R$.
\begin{proof}[Proof of Proposition~\ref{prop:mechanism-redundancy}]
The constant-row-sum assumption implies
$(K+\lambda I_M)\mathbf1_M=(\kappa+\lambda)\mathbf1_M$,
so $C\mathbf1_M=a\mathbf1_M$. Since $K$ is positive semidefinite,
$\kappa\geq0$ and hence $a>0$. Invariance of $\mathcal V$ means each
column of $CQ$ can be expressed using $Q$ and $\mathbf1_M$, which gives
$B$ and $c$. Using symmetry of $C$ in
Equation~\eqref{eq:mechanism-normalized-maps},
\[
 A_K(z)=
 \frac{(CQ)^\top k(z)}{(C\mathbf1_M)^\top k(z)}
 =\frac{B^\top Q^\top k(z)+c\mathbf1_M^\top k(z)}
 {a\mathbf1_M^\top k(z)}
 =\frac{B^\top A_R(z)+c}{a}.
\]
Subtract $A_A$ and multiply by $g$ to obtain the displayed departure
identity. Since $gA_A=g\mu+P\{gP^\top(z-\mu)\}$, all its coordinates
belong to $\mathcal S_0$. Thus the KAHM-augmented span is contained in
the RBF-augmented span.

For the reverse inclusion, the restriction of the invertible matrix $C$
to the invariant finite-dimensional space $\mathcal V$ is injective and
therefore surjective onto $\mathcal V$. Its images
$C\mathbf1_M=a\mathbf1_M$ and the columns of $CQ$ consequently span
$\mathcal V$. Hence there exist $D\in\R^{d\times d}$ and $e\in\R^d$ such that
$Q=(CQ)D+\mathbf1_M e^\top$. Multiplying its transpose by
$k(z)/(\mathbf1_M^\top k(z))$ gives
$A_R(z)=aD^\top A_K(z)+e$. After multiplication by $g$ and subtraction
of $gA_A$, this expresses each RBF departure coordinate using KAHM
departure coordinates and $\mathcal S_0$. The spans therefore agree,
including when the original reference-coordinate columns are redundant.
\end{proof}

For example, take scalar references $q=(-1,1)^\top$, $P=1$, $g=1$,
and $A_A(z)=z$. If the off-diagonal Gram entry is $\rho$, then
$K\mathbf1_2=(1+\rho)\mathbf1_2$ and $Kq=(1-\rho)q$. Therefore
\begin{equation}
 A_K=c_\rho A_R,\qquad
 d_K=c_\rho d_R+(c_\rho-1)z,\qquad
 c_\rho=\frac{1+\rho+\lambda}{1-\rho+\lambda}.
 \label{eq:mechanism-two-reference}
\end{equation}
The coefficient $c_\rho$ is positive, so
$\operatorname{span}(1,z,d_K)=\operatorname{span}(1,z,d_R)$.
With the sample-variance precision $1/2$, $\rho=e^{-1}$.
The regularization iteration used by the implementation yields
$\lambda=2.6519665429866075$ and $c_\rho=1.224037566462905$.
The maximum numerical identity error on the constructed query grid is
below $6\times10^{-16}$.

These are equalities between untransformed feature spans. An isotropic
ridge penalty is not invariant under a general nonorthogonal change of
coordinates. Applying separate nonlinear radial transformations can also
remove the displayed linear relation. A forecasting comparison therefore
evaluates the complete reconstruction, feature transformation, and fitting
procedure, even when two untransformed feature sets have the same span.

\subsection{An explicit distinct reconstruction basis}

\paragraph{Irregular scalar references.}
Let $q_1<q_2<q_3$ with $q_1+q_3\ne2q_2$, and define
$k_i(z)=\exp\{-t(z-q_i)^2/2\}$ for precision $t>0$.
Let $K_{ij}=\exp\{-t(q_i-q_j)^2/2\}$,
$C=(K+\lambda I_3)^{-1}$ for $\lambda>0$, and suppose every entry of
$v=C\mathbf1_3$ is positive. For the reconstructions in
Equation~\eqref{eq:mechanism-normalized-maps}, with $Q=q=(q_1,q_2,q_3)^\top$,
gate $g=1$, and $A_A(z)=z$,
\[
 A_K\notin\operatorname{span}(1,z,A_R),\qquad
 d_K\notin\operatorname{span}(1,z,d_R).
\]
These are statements about scalar functions on the whole real line.
\begin{proof}[Proof of Proposition~\ref{prop:mechanism-distinct}]
Set $u=Cq$. Since $C$ is symmetric,
$A_K(z)=(k(z)^\top u)/(k(z)^\top v)$, and positivity of $v$ makes
its denominator positive for every $z$. The six sums $q_i+q_j$ with
$i\leq j$ are distinct: strict ordering of the three references orders
all the sums except possibly $q_1+q_3$ and $2q_2$, and their equality
is excluded by assumption. Products of Gaussian translates satisfy
\[
 k_i(z)k_j(z)=
 e^{-tx^2}e^{-t(q_i^2+q_j^2)/2}e^{t(q_i+q_j)z}.
\]
After removal of the common positive factor $e^{-tx^2}$, these products
are nonzero multiples of exponentials with six distinct real exponents.
Such exponentials are linearly independent. Indeed, differentiating a
vanishing linear combination zero through fifth order at $z=0$ gives
a Vandermonde system whose determinant is the product of the nonzero
pairwise exponent differences.

Suppose first that $A_K=cA_R+d$ for scalars $c,d$. Define
$w=u-dv$. Multiplication by the two positive denominators gives
\[
 \left(\sum_i k_iw_i\right)\left(\sum_j k_j\right)
 =c\left(\sum_i k_iq_i\right)\left(\sum_j k_jv_j\right).
\]
Independence of the six products implies equality of their coefficients.
The diagonal coefficients give $w_i=cq_iv_i$, and the off-diagonal
coefficients give
$w_i+w_j=c(q_iv_j+q_jv_i)$ for $i<j$.
If $c=0$, all $w_i$ vanish, so $Cq=dC\mathbf1_3$. Invertibility of
$C$ would imply $q=d\mathbf1_3$, contrary to the distinct references.
If $c\ne0$, substituting the diagonal equations into the off-diagonal
ones yields $(q_i-q_j)(v_i-v_j)=0$. Thus $v=a\mathbf1_3$ for some
$a>0$, and $(K+\lambda I_3)v=\mathbf1_3$ implies that $K$ has
constant row sums. This is impossible: its middle-row sum exceeds the
first-row sum by $K_{23}-K_{13}>0$ and the last-row sum by
$K_{12}-K_{13}>0$, since the Gaussian decreases strictly with positive
distance. This rules out every affine relation $A_K=cA_R+d$.

Both reconstructions are bounded functions of $z$. The RBF reconstruction
is a convex combination of $q_1,q_2,q_3$. Also,
\[
 A_K(z)=\frac{\sum_i k_i(z)v_i\,(u_i/v_i)}
 {\sum_i k_i(z)v_i}
\]
is a convex combination of the three finite numbers $u_i/v_i$.
If $A_K=cA_R+d+bx$ held on the real line, boundedness of both
reconstructions would force $b=0$. The affine relation already ruled out
would then follow. This proves the first span statement. Finally,
$d_K=A_K-z$ and $d_R=A_R-z$ imply
$\operatorname{span}(1,z,d_R)=\operatorname{span}(1,z,A_R)$.
Membership of $d_K$ in this span would imply membership of $A_K$, a
contradiction.
\end{proof}

A concrete example uses $q=(-1,0,2)^\top$, sample-variance precision
$t=3/7$, and $\lambda=4.498274743339644$ from the specified
regularization iteration. Positivity of $C\mathbf1_3$ can be verified
without examining a query grid. Write $K=I_3+E$ and $a=1+\lambda$.
For a matrix, let $\|\cdot\|_\infty$ denote the maximum absolute row
sum; for a vector, it denotes the largest absolute entry. If
$a>2\|E\|_\infty$, then the Neumann series for
$aC=(I_3+E/a)^{-1}$ gives
\[
 \|aC\mathbf1_3-\mathbf1_3\|_\infty
 \leq\sum_{j=1}^{\infty}(\|E\|_\infty/a)^j
 =\frac{\|E\|_\infty/a}{1-\|E\|_\infty/a}<1.
\]
Every entry of $aC\mathbf1_3$ is then strictly positive. In the example,
$\|E\|_\infty=e^{-3/14}+e^{-6/7}=1.2314905926823394$, so the
condition holds. On $81$ constructed queries in $[-3,4]$, the maximum
absolute residual after least-squares projection of the sampled $d_K$
values onto the sampled $(1,z,d_R)$ columns is
$0.05355670997078987$. This numerical illustration supplements the
analytic existence result. It does not establish distinctness for every
higher-dimensional bank used in the forecasting experiments.

\subsection{A bounded correction with a free intercept}

For a feature block $v\in\R^{p_b}$ of positive width $p_b$, define
\begin{equation}
 T_{R_b}(v)=\frac{v}{\sqrt{1+\|v\|_2^2/R_b^2}},\qquad
 R_b=\sqrt{p_b}.
 \label{eq:mechanism-radial}
\end{equation}
Its norm is at most $R_b$. To see its sensitivity to its own input, set
$a=1+\|v\|_2^2/R_b^2$. Its Jacobian is
$a^{-1/2}I-a^{-3/2}vv^\top/R_b^2$. For $v\ne0$, the eigenvalue in
the direction of $v$ is $a^{-3/2}$, and each orthogonal direction has
eigenvalue $a^{-1/2}$; at zero, the Jacobian is the identity. All its
singular values are therefore at most one. Integrating the derivative
along the segment between any two inputs proves
$\|T_{R_b}(v)-T_{R_b}(v')\|_2\leq\|v-v'\|_2$.
Thus the radial transformation is nonexpansive.

Concatenating separately transformed blocks gives a feature vector
$\psi(z)\in\R^p$, where $p=\sum_b p_b$, satisfying
$\|\psi(z)\|_2\leq R=(\sum_bR_b^2)^{1/2}$. This norm bound holds
wherever the preceding feature construction is defined. Nonexpansiveness
of the radial transformation does not by itself give a Lipschitz constant
for the full history-to-feature map, which includes the signed
reconstruction denominator and its fallback rule.

\paragraph{Centered ridge correction.}
For $n\geq1$, let $e_i\in\R$ be residual targets and
$\psi_i\in\R^p$ be training feature vectors. Fit $b\in\R$ and
$w\in\R^p$ by minimizing
\[
 \frac1n\sum_{i=1}^n(e_i-b-\psi_i^\top w)^2+\alpha\|w\|_2^2,
 \qquad\alpha>0,
\]
with no penalty on $b$. Define the training means
$\bar e=n^{-1}\sum_i e_i$ and
$\bar\psi=n^{-1}\sum_i\psi_i$, the residual vector
$e=(e_1,\ldots,e_n)^\top$, its centered version
$e_c=e-\bar e\mathbf1_n$, and its centered root mean square
$\sigma_e=\|e_c\|_2/\sqrt n$. Suppose query features satisfy
$\|\psi(z)\|_2\leq R$. Then, for any scalar $\gamma$,
\begin{equation}
 \|w\|_2\leq\frac{\sigma_e}{2\sqrt\alpha},\qquad
 |\gamma q(z)|\leq|\gamma|\left[
 |\bar e|+(R+\|\bar\psi\|_2)\frac{\sigma_e}{2\sqrt\alpha}\right],
 \quad q(z)=\bar e+(\psi(z)-\bar\psi)^\top w.
 \label{eq:mechanism-correction-bound}
\end{equation}
\begin{proof}[Proof of Proposition~\ref{prop:mechanism-correction-bound}]
For fixed $w$, differentiating with respect to the intercept gives
$b=\bar e-\bar\psi^\top w$. Let $\Psi_c\in\R^{n\times p}$ have
row $i$ equal to $(\psi_i-\bar\psi)^\top$. Substituting the fitted
intercept reduces the objective to
$n^{-1}\|e_c-\Psi_cw\|_2^2+\alpha\|w\|_2^2$.
Its unique minimizer is
\[
 w=(\Psi_c^\top\Psi_c+n\alpha I_p)^{-1}\Psi_c^\top e_c.
\]
For a singular value $s\geq0$ of $\Psi_c$, the corresponding singular
value of the coefficient operator is $s/(s^2+n\alpha)$. The inequality
$(s-\sqrt{n\alpha})^2\geq0$ implies
\[
 \frac{s}{s^2+n\alpha}\leq\frac{1}{2\sqrt{n\alpha}}.
\]
The operator norm is therefore at most $1/(2\sqrt{n\alpha})$, giving
$\|w\|_2\leq\|e_c\|_2/(2\sqrt{n\alpha})
=\sigma_e/(2\sqrt\alpha)$. Finally, the fitted intercept gives
$q(z)=\bar e+(\psi(z)-\bar\psi)^\top w$.
Cauchy--Schwarz and the triangle inequality yield
\[
 |q(z)|\leq|\bar e|+
 (\|\psi(z)\|_2+\|\bar\psi\|_2)\|w\|_2
 \leq|\bar e|+(R+\|\bar\psi\|_2)
 \frac{\sigma_e}{2\sqrt\alpha}.
\]
Multiplication by $|\gamma|$ proves the correction bound.
\end{proof}

The mean-squared-error objective corresponds to a summed-squared-error
objective with penalty $n\alpha$. Keeping the numerical value $\alpha$
unchanged when switching to a summed loss would define a different
estimator. For multiple forecast horizons, the bound applies separately
to each fitted residual vector and each selected pair
$(\alpha_h,\gamma_h)$. This is a general property of ridge regression
on bounded features, including the KAHM features considered here. It
bounds the additive correction but does not bound the retained raw
baseline or total forecast error. The experiment checks the feature,
coefficient, and correction inequalities on every saved task and reports
forecast errors separately.
